
\documentclass{article}
\usepackage{colortbl}
\usepackage{makecell}
\usepackage{multirow}
\usepackage{supertabular}

\begin{document}

\newcounter{utterance}

\centering \large Interaction Transcript for game `cladder', experiment `full\_v1.5\_default', episode 6776 with Qwen3.5{-}9B.
\vspace{24pt}

{ \footnotesize  \setcounter{utterance}{1}
\setlength{\tabcolsep}{0pt}
\begin{supertabular}{c@{$\;$}|p{.15\linewidth}@{}p{.15\linewidth}p{.15\linewidth}p{.15\linewidth}p{.15\linewidth}p{.15\linewidth}}
    \# & \multicolumn{2}{c}{Player} && \multicolumn{2}{c}{Game Master} \\
    \hline

    \theutterance \stepcounter{utterance}  
    & & & \multicolumn{4}{p{0.6\linewidth}}{
        \cellcolor[rgb]{0.9,0.9,0.9}{
            \makecell[{{p{\linewidth}}}]{
                \texttt{\tiny{[P1$\langle$GM]}}
                \texttt{You are an expert in causal inference. The following question is not a typical commonsense query, but rather a meticulously designed question created by a professor specializing in causal inference, intended to assess the students' mastery of the course content.} \\
\texttt{Imagine a self{-}contained, hypothetical world with only the following conditions, and without any unmentioned factors or causal relationships: Gwet has a direct effect on yomx and xevu. Yomx has a direct effect on xevu. For those who are not gwet and are not yomx, the probability of xevu is 82\%. For those who are not gwet and are yomx, the probability of xevu is 63\%. For those who are gwet and are not yomx, the probability of xevu is 36\%. For those who are gwet and are yomx, the probability of xevu is 18\%. The overall probability of gwet is 41\%. Will yomx increase the chance of xevu?} \\
\texttt{Start your answer with 'Yes' or 'No', followed by additional reasoning or evidence to support your explanation.} \\
            }
        }
    }
     \\ \\

    \theutterance \stepcounter{utterance}  
    & \multicolumn{4}{p{0.6\linewidth}}{
        \cellcolor[rgb]{0.561,0.737,0.561}{
            \makecell[{{p{\linewidth}}}]{
                \texttt{\tiny{[P1$\rangle$GM]}}
                \texttt{No} \\
\\ 
\texttt{**Reasoning:**} \\
\\ 
\texttt{To determine if \$yomx\$ increases the chance of \$xevu\$, we must calculate the causal effect of \$yomx\$ on \$xevu\$. In this hypothetical world, \$gwet\$ is a confounder because it has a direct effect on both \$yomx\$ (implied by the stratification logic where \$gwet\$ status defines the groups) and \$xevu\$. More critically, the problem states \$gwet\$ has a direct effect on \$yomx\$ and \$yomx\$ has a direct effect on \$xevu\$. To isolate the causal effect of \$yomx\$ on \$xevu\$, we must compare the probability of \$xevu\$ given \$yomx=1\$ versus \$yomx=0\$ while holding the confounder \$gwet\$ constant (backdoor adjustment).} \\
\\ 
\texttt{Let \$P(xevu | gwet, yomx)\$ denote the conditional probabilities provided:} \\
\texttt{1.  \$P(xevu | \textbackslash{}neg gwet, \textbackslash{}neg yomx) = 0.82\$} \\
\texttt{2.  \$P(xevu | \textbackslash{}neg gwet, yomx) = 0.63\$} \\
\texttt{3.  \$P(xevu | gwet, \textbackslash{}neg yomx) = 0.36\$} \\
\texttt{4.  \$P(xevu | gwet} \\
            }
        }
    }
    & & \\ \\

    \theutterance \stepcounter{utterance}  
    & & & \multicolumn{2}{p{0.3\linewidth}}{
        \cellcolor[rgb]{0.9,0.9,0.9}{
            \makecell[{{p{\linewidth}}}]{
                \texttt{\tiny{[GM$|$GM]}}
                \texttt{No} \\
            }
        }
    }
    & & \\ \\

    \theutterance \stepcounter{utterance}  
    & & & \multicolumn{2}{p{0.3\linewidth}}{
        \cellcolor[rgb]{0.9,0.9,0.9}{
            \makecell[{{p{\linewidth}}}]{
                \texttt{\tiny{[GM$|$GM]}}
                \texttt{no} \\
            }
        }
    }
    & & \\ \\

    \theutterance \stepcounter{utterance}  
    & & & \multicolumn{2}{p{0.3\linewidth}}{
        \cellcolor[rgb]{0.9,0.9,0.9}{
            \makecell[{{p{\linewidth}}}]{
                \texttt{\tiny{[GM$|$GM]}}
                \texttt{game\_result = WIN} \\
            }
        }
    }
    & & \\ \\

\end{supertabular}
}

\end{document}
